% ============================================================================
%  3. DIAGRAMA ARQUITECTÓNICO FINAL
%  ----------------------------------------------------------------------------
%  Fuente: TSO-II-final/docs/a10.md, ajustado con los datos de red
%  vigentes del proyecto (AGENTS.md).
% ============================================================================

\section{Diagrama arquitectónico final}

La infraestructura del banco Sudoers se diseñó sobre una red local única
(\texttt{192.168.0.0/24}) que se segmenta l\'ogicamente por roles y nivel de
acceso, de modo que cada zona expone \'unicamente los servicios que
necesita. Esta secci\'on documenta la segmentaci\'on adoptada, los roles
asignados a cada servidor, el flujo de comunicaci\'on resultante y la
justificaci\'on t\'ecnica del dise\~no.

\subsection{Segmentaci\'on de red}

La red se divide en cuatro zonas funcionales, delimitadas por rangos de
direcciones y pol\'iticas de filtrado distintas, tal como se detalla en la
tabla \ref{tab:zonas-red}.

\begin{table}[H]
\centering
\caption{Segmentaci\'on de la red local por zonas funcionales}
\label{tab:zonas-red}
\contenidoConFuente{%
  \begin{tablaAPA}
    \begin{tabular}{@{}p{2.4cm}p{4.2cm}p{3.4cm}p{4cm}@{}}
    \toprule
    \textbf{Zona} & \textbf{Rango} & \textbf{Funci\'on} & \textbf{Equipos} \\
    \midrule
    A & \texttt{192.168.0.1} & Per\'imeter y salida a Internet & Router TP-Link TL-WR850N \\
    B & \texttt{192.168.0.2} & Host servidor de dominio y servicios & \texttt{dc1} (Debian 12) \\
    B & \texttt{192.168.0.3} & Monitoreo de infraestructura & \texttt{zabbix} (Debian 13) \\
    C & \texttt{.60}--\texttt{.99} & Puestos operativos y cajas & Terminales Linux Kiosk \\
    D & \texttt{.50}--\texttt{.52} & Administraci\'on de TI & \texttt{pc-joaquin}, \texttt{pc-david}, \texttt{pc-nicolas} \\
    D & \texttt{.100}--\texttt{.199} & Administraci\'on y gerencia (din\'amico) & Estaciones Windows 10/11 \\
    \bottomrule
    \end{tabular}
  \end{tablaAPA}%
}{Elaboraci\'on propia.}
\end{table}

El servicio DHCP del router (Zona A) se desactiva de forma deliberada para
que el servidor Kea sea el \'unico emisor de ofertas de la red, evitando as\'\i
conflicto de doble respuesta \textsc{Discover} \cite{rfc2131}.

\subsection{Servidores y roles}

La tabla \ref{tab:servidores} detalla cada servidor, su direcci\'on y el rol
que ejerce dentro de la arquitectura.

\begin{table}[H]
\centering
\caption{Servidores, direcciones y roles en la arquitectura final}
\label{tab:servidores}
\contenidoConFuente{%
  \begin{tablaAPA}
    \begin{tabular}{@{}p{2.2cm}p{3.2cm}p{2.4cm}p{5.8cm}@{}}
    \toprule
    \textbf{Host} & \textbf{IP} & \textbf{SO} & \textbf{Rol} \\
    \midrule
    \texttt{dc1} & \texttt{192.168.0.2} & Debian 12 & Controlador de dominio Samba 4, DNS, NTP, firewall, servicios en Docker \\
    \texttt{zabbix} & \texttt{192.168.0.3} & Debian 13 & Servidor de monitoreo Zabbix (server, web y PostgreSQL) \\
    Router & \texttt{192.168.0.1} & Firmware & Puerta de enlace, NAT a Internet y segmentaci\'on \\
    \bottomrule
    \end{tabular}
  \end{tablaAPA}%
}{Elaboraci\'on propia.}
\end{table}

El host \texttt{dc1} concentra una arquitectura de tres capas de
procesamiento. La capa nativa ejecuta el controlador de dominio Samba 4, que
integra en un \'unico daemon los servicios de LDAP, Kerberos, DNS y SMB
\cite{sambateam2024}, junto con \texttt{chrony} para la sincronizaci\'on
horaria --imprescindible para Kerberos, que rechaza los tickets con un
desfase de reloj superior a cinco minutos--, el firewall \texttt{nftables} y el servicio
\texttt{sshd} endurecido. La segunda capa corresponde a los contenedores
Docker, donde residen el proxy inverso Nginx, el servidor DHCP Kea en modo
de red \texttt{host}, la base de datos PostgreSQL, el servicio de correo
Postfix y Dovecot, el gestor de impresi\'on CUPS y el panel de gesti\'on
Portainer. La tercera capa constituye el servidor \texttt{zabbix}, que
opera como m\'aquina independiente y es alcanzable \'unicamente a trav\'es
del proxy inverso.

La decisi\'on de ejecutar el controlador de dominio de forma nativa, y no
contenedorizado, responde a la fragilidad de Samba frente al aislamiento de
redes y sistemas de archivos en contenedores \cite{sambateam2024}.

\subsection{Flujo de comunicaci\'on entre servicios}

La figura \ref{fig:arquitectura-zonas} representa la topolog\'ia general de
la red y las relaciones entre el per\'imeter, el host de servicios, el
servidor de monitoreo y las zonas de puestos de trabajo.

\capturaAPA{fig:arquitectura-zonas}{Diagrama de la arquitectura de red segmentada por zonas y roles de acceso}{captura-arquitectura-red.jpeg}{Diagrama de la topolog\'ia general de la red del banco Sudoers, con la segmentaci\'on en zonas A (per\'imeter), B (servidores \texttt{dc1} y \texttt{zabbix}), C (puestos de cajas) y D (administraci\'on y gerencia), junto con los flujos de comunicaci\'on entre ellas.}

La matriz de la tabla \ref{tab:matriz-acceso} formaliza el flujo anterior:
para cada servicio se indican el puerto y protocolo expuestos, la capa de
ejecici\'on y el nivel de acceso permitido.

\begin{table}[H]
\centering
\caption{Matriz de acceso y visibilidad de los servicios}
\label{tab:matriz-acceso}
\contenidoConFuente{%
  \begin{tablaAPA}
    \begin{tabular}{@{}p{4.9cm}p{2.4cm}p{2.1cm}p{4.5cm}@{}}
    \toprule
    \textbf{Servicio} & \textbf{Puerto} & \textbf{Capa} & \textbf{Acceso} \\
    \midrule
    Samba DNS / LDAP / Kerberos & \texttt{53}, \texttt{88}, \texttt{389} & Nativa & Todo el dominio (validaci\'on) \\
    Recursos compartidos SMB & \texttt{445/tcp} & Nativa & Filtrado por grupos de usuario AD \\
    SSH / nftables & \texttt{22/tcp} & Nativa & Exclusivo de administradores TI (\texttt{.50}--\texttt{.52}) \\
    Nginx proxy inverso & \texttt{80}, \texttt{443/tcp} & Docker & LAN operativa y gerencia \\
    Postfix / Dovecot & \texttt{25}, \texttt{143}, \texttt{587} & Docker & Usuarios autenticados del dominio \\
    Kea DHCP & \texttt{67/udp} & Docker (host) & Difusi\'on de toda la LAN \\
    PostgreSQL & \texttt{5432/tcp} & Docker & Aislado (solo Core Bancario) \\
    CUPS & \texttt{631/tcp} & Docker & Puestos de cajas y administraci\'on \\
    Portainer Web & \texttt{9443/tcp} & Docker & Exclusivo de administradores TI \\
    Zabbix (agente pasivo) & \texttt{10050/tcp} & Servidor \texttt{zabbix} & Clientes y host \\
    \bottomrule
    \end{tabular}
  \end{tablaAPA}%
}{Elaboraci\'on propia.}
\end{table}

\subsection{Justificaci\'on t\'ecnica del dise\~no}

El dise\~no se apoya en cinco decisiones t\'ecnicas fundamentadas.

\begin{enumerate}
  \item \textbf{Segmentaci\'on por roles.} Dividir la red en zonas con
        pol\'iticas de filtrado distintas reduce la superficie de ataque:
        los puestos de cajas no requieren acceso al panel de administraci\'on
        de contenedores ni a la consola del servidor de monitoreo.
  \item \textbf{Seguridad por defecto.} La pol\'itica \textsc{drop} por
        defecto de \texttt{nftables} y la restricci\'n del puerto
        \texttt{22/tcp} \'unicamente a los rangos de administraci\'on
        implementan el principio de m\'inimo privilegio
        \cite{stewart2017}.
  \item \textbf{Identidad centralizada.} Delegar la autenticaci\'on en
        Samba 4 AD DC sobre Kerberos y LDAP permite una gesti\'on \'unica de
        usuarios, grupos y permisos, en lugar de credenciales dispersas por
        cada servicio \cite{rfc4120, rfc4511}.
  \item \textbf{Aislamiento de datos sensibles.} PostgreSQL queda
        confinado a la red interna de Docker y el servicio de monitoreo se
        publica \'unicamente a trav\'es del proxy inverso, de modo que
        ninguno es alcanzable de forma directa desde la LAN.
  \item \textbf{Un \'unico emisor de direcciones.} Desactivar el DHCP del
        router y centralizar la asignaci\'on en Kea elimina la ambig\"uedad
        en la obtenci\'on de la direcci\'on, puerta de enlace y servidor DNS
        por parte de los clientes \cite{rfc2131, keaDocs}.
\end{enumerate}

Estas decisiones culminan, en conjunto, en una arquitectura que puede escalar
con nuevos servicios sobre la misma base de segmentaci\'on y controles de
acceso, principios que se verifican de forma pr\'actica en el cap\'itulo de
desarrollo t\'ecnico.
